% =====================================================================
%  Developer Guide: technical reference for the code base.
%  Build: ./scripts/build_docs.sh   (or compile in docs/ with pdflatex)
% =====================================================================
\documentclass[11pt,a4paper]{article}
\usepackage{campusdoc}

% Function reference entry:  \fn{signature}{description}
\newcommand{\fn}[2]{\par\smallskip\noindent\colorbox{codebg}{\parbox{\dimexpr\linewidth-2\fboxsep}{\ttfamily\small\detokenize{#1}}}\par\nopagebreak\smallskip
  \noindent\begin{minipage}{\linewidth}\setlength{\parindent}{0pt}\leftskip=1.2em #2\end{minipage}\par\medskip}

\begin{document}

\makedoctitle{Developer Guide}{Architecture, code reference and conventions}{CEMP-DG-01}

\tableofcontents
\newpage

% =====================================================================
\section{Introduction}

This guide is the technical reference for the Campus Event Management Platform. It explains how
the code is organised, what every module, route and function does, how the database is designed,
and how to extend the application safely. It is written for someone reading or changing the code.
How to \emph{use} the site is in the User Manual; how to install it is in the Installation \&
Operations Guide; the tests are documented in the Test Plan \& Report.

\subsection{Design principles}

\begin{enumerate}
  \item \textbf{Simple and explainable.} Plain functions, short modules, comments that say
        \emph{why}. No clever abstractions; every line should be explainable in a viva.
  \item \textbf{Raw SQL, no ORM.} All database work uses Python's built-in \code{sqlite3} module
        with hand-written SQL, so the queries can be read, tested, profiled and \code{EXPLAIN}ed.
  \item \textbf{Every value through a placeholder.} User input reaches SQLite only through
        \code{?} parameters, never through string formatting.
  \item \textbf{The database enforces the rules too.} \code{UNIQUE}, \code{CHECK} and foreign-key
        constraints back up the Python validation.
  \item \textbf{Hand-written HTML and CSS.} Jinja2 templates and one stylesheet; no CSS or
        JavaScript framework.
  \item \textbf{Measure before optimising.} Performance changes are justified by
        \code{cProfile} measurements and query plans kept in \code{profiling/}.
\end{enumerate}

\subsection{Technology stack}

\begin{table}[H]
\centering
\small
\begin{tabularx}{\linewidth}{@{}l X@{}}
\toprule
\textbf{Layer} & \textbf{Technology} \\
\midrule
Language        & Python 3.10+ \\
Web framework   & Flask 3 (routing, Blueprints, Jinja2 templates, Click CLI) \\
Authentication  & Flask-Login (sessions), Werkzeug \code{generate_password_hash} / \code{check_password_hash} \\
Database        & SQLite through \code{sqlite3}, raw SQL \\
Front end       & Jinja2 HTML templates, one hand-written CSS file \\
Testing         & pytest with Flask's test client \\
Demo data       & Faker (\code{en_IN} locale, fixed seed) \\
Profiling       & \code{cProfile}, \code{pstats}, SnakeViz, \code{EXPLAIN QUERY PLAN} \\
Automation      & Bash scripts \\
Documentation   & LaTeX (article and Beamer), Markdown \\
\bottomrule
\end{tabularx}
\caption{Technology stack.}
\end{table}

% =====================================================================
\section{Repository layout}

\begin{lstlisting}[language={}, caption={Project tree (generated files such as venv/ and instance/ are git-ignored).}]
Campus-Event-Management-Platform/
|-- app/                      Flask application package
|   |-- __init__.py           create_app(): config, blueprints, error pages, home page
|   |-- db.py                 get_db/close_db, init-db, add-indexes, create-admin commands
|   |-- auth.py               User class, login manager, admin_required, register/login/logout
|   |-- events.py             search, event detail, register/cancel, my registrations
|   |-- admin.py              dashboard, venue + event CRUD, clash check, participants, CSV
|   |-- reports.py            report queries (JOIN + GROUP BY)
|   |-- seed.py               "flask seed" demo-data command
|   |-- templates/            Jinja2 templates: base, index, error, auth/, events/, admin/
|   `-- static/style.css      the only stylesheet
|-- sql/
|   |-- schema.sql            tables and constraints
|   `-- indexes.sql           performance indexes (added after profiling)
|-- scripts/                  Bash automation: common, setup, run, seed, test, profile, build_docs
|-- tests/                    conftest.py (fixtures) + 5 test files, 45 tests
|-- profiling/                run_profile.py, explain.py, RESULTS.md, before/ rewrite_only/ after/
|-- report/                   report.tex, slides.tex, figures/
|-- docs/                     campusdoc.sty + the four guides (this file is one of them)
|-- instance/                 campus.db (created by setup.sh, git-ignored)
|-- requirements.txt  pytest.ini  .gitignore  .gitattributes
`-- README.md  AI_USAGE.md  CLAUDE.md
\end{lstlisting}

% =====================================================================
\section{Architecture}

\subsection{Overview}

The application follows the usual three-tier structure of a server-rendered web app:
\textbf{presentation} (Jinja2 templates and CSS), \textbf{application logic} (Flask Blueprints and
plain Python functions) and \textbf{data} (SQLite). There is no client-side JavaScript apart from the
browser's built-in \code{confirm()} on delete buttons.

\begin{figure}[H]
\centering
\resizebox{\linewidth}{!}{%
\begin{tikzpicture}[
    font=\small,
    box/.style={draw, rounded corners=3pt, fill=primary!6, align=center, minimum height=0.9cm, text width=2.4cm},
    mod/.style={draw, rounded corners=2pt, fill=white, align=center, minimum height=0.8cm, text width=2.7cm},
    arr/.style={-{Stealth[length=2.5mm]}, thick},
  ]
  \node[box, minimum height=1.6cm] (browser) at (-0.8,-0.4) {\textbf{Browser}\\ HTML forms\\ (GET / POST)};
  \node[mod] (init)   at (4.6,2.2)  {\code{__init__.py}\\ \code{create_app()}};
  \node[mod] (auth)   at (4.6,0.9)  {\code{auth} Blueprint\\ \route{/auth/...}};
  \node[mod] (events) at (4.6,-0.4) {\code{events} Blueprint\\ \route{/events/...}};
  \node[mod] (admin)  at (4.6,-1.7) {\code{admin} Blueprint\\ \route{/admin/...}};
  \node[mod] (reports) at (8.3,-1.7) {\code{reports.py}};
  \node[mod] (db)     at (8.3,0.25) {\code{db.py}\\ \code{get_db()}};
  \node[mod] (tpl)    at (8.3,2.2)  {Jinja2 templates\\ + \code{style.css}};
  \node[draw, dashed, rounded corners=5pt, inner sep=7pt, fit=(init)(admin)(reports)(tpl),
        label={[font=\small\bfseries]above:Flask application}] {};
  \node[draw, shape=cylinder, shape border rotate=90, aspect=0.25, fill=orange!15,
        minimum height=1.8cm, minimum width=1.9cm, align=center] (sqlite) at (13.0,0.25)
        {\textbf{SQLite}\\ \code{campus.db}};
  \draw[arr] ([yshift=4pt]browser.east) -- node[above, font=\footnotesize, sloped]{request} ([yshift=4pt]events.west);
  \draw[arr] ([yshift=-4pt]events.west) -- node[below, font=\footnotesize, sloped]{HTML page} ([yshift=-4pt]browser.east);
  \draw[arr] (auth.east) -- (db.west);
  \draw[arr] (events.east) -- (db.west);
  \draw[arr] (admin.east) -- (reports.west);
  \draw[arr] (admin.north east) -- (db.south west);
  \draw[arr] (reports.north) -- (db.south);
  \draw[arr] (db.east) -- node[above, font=\footnotesize]{SQL + \code{?}} (sqlite.west);
\end{tikzpicture}}
\caption{Components and how a request flows through them.}
\end{figure}

\subsection{Life of a request}

Taking ``a student clicks \textsf{Register}'' as an example:

\begin{enumerate}
  \item The browser sends \code{POST /events/7/register} with the session cookie.
  \item Flask matches the URL rule \route{/events/<int:event_id>/register}, so it calls the
        view function \code{register_for_event} in \code{events.py} with \code{event_id=7}.
  \item \code{@login_required} asks Flask-Login for the current user. Flask-Login reads the user id
        from the signed cookie and calls \code{load_user()}, which runs one \code{SELECT} on
        \code{users}. A visitor is redirected to \route{/auth/login}.
  \item The view calls \code{get_event(7)}. \code{get_db()} opens a connection the first time it is
        needed in this request and stores it in \code{flask.g}.
  \item The view checks the registration rules (Section~\ref{sec:rules}), runs the
        \code{INSERT}, commits, and \code{flash()}es a message.
  \item It returns a redirect to the detail page (\emph{Post/Redirect/Get}), so refreshing the page
        never submits the form twice.
  \item When the request ends, \code{close_db()} (registered with \code{teardown_appcontext})
        closes the connection.
  \item The browser follows the redirect; the detail page renders the flashed message from
        \code{base.html}.
\end{enumerate}

\subsection{Application factory}

\code{create_app(test_config=None)} in \code{app/__init__.py} builds a new Flask application:

\begin{enumerate}
  \item creates \code{Flask(__name__, instance_relative_config=True)}, so \code{app.instance_path}
        is the \code{instance/} folder next to \code{app/};
  \item sets \code{SECRET_KEY} (env var, default \code{dev-change-me}), \code{DATABASE}
        (env var \code{CAMPUS_DB}, default \code{instance/campus.db}) and
        \code{SESSION_COOKIE_SAMESITE="Lax"};
  \item applies \code{test_config} on top (the tests pass a temporary database path);
  \item creates the \code{instance/} folder;
  \item registers the database helpers and CLI commands (\code{db.init_app}) and the \code{seed}
        command;
  \item initialises Flask-Login and registers the \code{auth}, \code{events} and \code{admin}
        Blueprints;
  \item defines the home page route \route{/} and the 403 and 404 error handlers, which render
        \code{error.html}.
\end{enumerate}

Building the app in a function instead of at import time is what lets every test get its own
isolated app and database.

% =====================================================================
\section{Database design}

\subsection{Entity--relationship model}

\begin{figure}[H]
\centering
\begin{tikzpicture}[
    entity/.style={draw, rounded corners=2pt, fill=primary!4, font=\small, inner sep=2pt},
    rel/.style={draw, thick},
    card/.style={font=\footnotesize\bfseries, inner sep=1pt},
    lbl/.style={font=\footnotesize\itshape, fill=white, inner sep=1pt},
  ]
  \node[entity] (users) at (0,0) {\begin{tabular}{l}\textbf{users}\\ \hline id (PK)\\ name\\ email (UNIQUE)\\ password\_hash\\ role\end{tabular}};
  \node[entity] (events) at (8.5,0) {\begin{tabular}{l}\textbf{events}\\ \hline id (PK)\\ title, description\\ category\\ venue\_id (FK)\\ start\_time, end\_time\\ capacity\\ created\_by (FK)\end{tabular}};
  \node[entity] (regs) at (0,-4.6) {\begin{tabular}{l}\textbf{registrations}\\ \hline id (PK)\\ user\_id (FK)\\ event\_id (FK)\\ registered\_at\end{tabular}};
  \node[entity] (venues) at (8.5,-4.6) {\begin{tabular}{l}\textbf{venues}\\ \hline id (PK)\\ name\\ capacity\\ location\end{tabular}};
  \draw[rel] (users) -- node[lbl]{creates} node[card, pos=0.06, above]{1} node[card, pos=0.94, above]{N} (events);
  \draw[rel] (users) -- node[lbl]{makes} node[card, pos=0.1, right]{1} node[card, pos=0.9, right]{N} (regs);
  \draw[rel] (regs) -- node[lbl]{for} node[card, pos=0.1, above]{N} node[card, pos=0.9, above]{1} (events);
  \draw[rel] (events) -- node[lbl]{held at} node[card, pos=0.1, right]{N} node[card, pos=0.9, right]{1} (venues);
\end{tikzpicture}
\caption{Four tables. \texttt{registrations} links students and events (many-to-many).}
\end{figure}

\subsection{Table reference}

The schema is in \code{sql/schema.sql}. Running it drops and recreates all tables, children first.

\newcommand{\coltable}[1]{%
\begin{tabularx}{\linewidth}{@{}>{\ttfamily}l >{\ttfamily}l X@{}}
\toprule
\textbf{\textrm{Column}} & \textbf{\textrm{Type}} & \textbf{Constraints and meaning} \\
\midrule
#1
\bottomrule
\end{tabularx}}

\paragraph{users}: everyone who can log in.\par\smallskip
\noindent\small\coltable{
id & INTEGER & Primary key, auto-increment. \\
name & TEXT & \code{NOT NULL}. Display name. \\
email & TEXT & \code{NOT NULL UNIQUE}. Stored lower-case by the app; the login name. \\
password\_hash & TEXT & \code{NOT NULL}. Werkzeug salted hash (scrypt by default); never the password. \\
role & TEXT & \code{NOT NULL DEFAULT 'student'}, \code{CHECK (role IN ('admin','student'))}. \\
}\normalsize

\paragraph{venues}: physical places.\par\smallskip
\noindent\small\coltable{
id & INTEGER & Primary key, auto-increment. \\
name & TEXT & \code{NOT NULL}. \\
capacity & INTEGER & \code{NOT NULL}, \code{CHECK (capacity > 0)}. \\
location & TEXT & \code{NOT NULL}. Building and floor. \\
}\normalsize

\paragraph{events}: things students register for.\par\smallskip
\noindent\small\coltable{
id & INTEGER & Primary key, auto-increment. \\
title & TEXT & \code{NOT NULL}. \\
description & TEXT & Optional. \\
category & TEXT & \code{NOT NULL}. One of \code{CATEGORIES} (checked in Python). \\
venue\_id & INTEGER & \code{NOT NULL}, foreign key to \code{venues(id)}. \\
start\_time & TEXT & \code{NOT NULL}. \code{'YYYY-MM-DD HH:MM'}. \\
end\_time & TEXT & \code{NOT NULL}. Same format. Table check \code{CHECK (end_time > start_time)}. \\
capacity & INTEGER & \code{NOT NULL}, \code{CHECK (capacity > 0)}. At most the venue's capacity (checked in Python). \\
created\_by & INTEGER & \code{NOT NULL}, foreign key to \code{users(id)}: the admin who created it. \\
}\normalsize

\paragraph{registrations}: which student is registered for which event.\par\smallskip
\noindent\small\coltable{
id & INTEGER & Primary key, auto-increment. \\
user\_id & INTEGER & \code{NOT NULL}, foreign key to \code{users(id)} \code{ON DELETE CASCADE}. \\
event\_id & INTEGER & \code{NOT NULL}, foreign key to \code{events(id)} \code{ON DELETE CASCADE}. \\
registered\_at & TEXT & \code{NOT NULL DEFAULT (datetime('now','localtime'))}. \\
\multicolumn{2}{@{}l}{\textrm{\emph{table constraint}}} & \code{UNIQUE (user_id, event_id)}: one registration per student per event. \\
}\normalsize

\subsection{Conventions}

\begin{description}
  \item[Times as text.] Times are \code{TEXT} in \code{'YYYY-MM-DD HH:MM'} format. Because the most
        significant part comes first, string comparison equals time comparison, so \code{<},
        \code{>}, \code{BETWEEN} and \code{ORDER BY} work directly. \code{now_string()} in
        \code{events.py} produces the current time in the same format. HTML
        \code{datetime-local} inputs use a \code{T} separator, converted by \code{to_db_time()}.
  \item[Foreign keys.] SQLite enforces them only when \code{PRAGMA foreign_keys = ON} is run on the
        connection; \code{get_db()} does this every time.
  \item[Rows as dicts.] \code{row_factory = sqlite3.Row}, so code writes \code{row["title"]}.
        Report functions convert rows with \code{dict(row)} before adding computed keys.
  \item[Seats left] are never stored; they are computed as \code{capacity - COUNT(r.id)} with a
        \code{LEFT JOIN} on \code{registrations}, so they can never get out of step.
\end{description}

\subsection{Indexes}

\code{sql/indexes.sql} is deliberately separate from the schema: the indexes were the result of
profiling, and keeping them apart lets the profiler measure with and without them
(\code{APPLY_INDEXES=0}).

\begin{table}[H]
\centering
\small
\begin{tabularx}{\linewidth}{@{}l l X@{}}
\toprule
\textbf{Index} & \textbf{Columns} & \textbf{Speeds up} \\
\midrule
\code{idx_registrations_event} & \code{registrations(event_id)} & every seats-left count and report \code{JOIN}; covering index for counts \\
\code{idx_events_venue_start}  & \code{events(venue_id, start_time)} & venue filter in search, venue clash check \\
\code{idx_events_category}     & \code{events(category)} & category filter in search \\
\code{sqlite_autoindex_registrations_1} & \code{(user_id, event_id)} & created automatically by the \code{UNIQUE} constraint; serves ``is this student registered?'' and My registrations \\
\bottomrule
\end{tabularx}
\caption{Indexes. An extra index on \texttt{registrations(user\_id)} was rejected as redundant.}
\end{table}

% =====================================================================
\section{Route reference}

{\small
\begin{xltabular}{\linewidth}{@{}>{\raggedright\arraybackslash}p{0.37\linewidth} l l X@{}}
\toprule
\textbf{URL} & \textbf{Method} & \textbf{Access} & \textbf{Endpoint and purpose} \\
\midrule
\endhead
\route{/}                                   & GET      & all     & \code{index}: home page with counts \\
\midrule
\route{/auth/register}                      & GET/POST & all     & \code{auth.register}: student sign-up \\
\route{/auth/login}                         & GET/POST & all     & \code{auth.login} \\
\route{/auth/logout}                        & GET      & login   & \code{auth.logout} \\
\midrule
\route{/events/}                            & GET      & all     & \code{events.list_events}: upcoming events; query params \code{q}, \code{category}, \code{venue_id}, \code{date_from}, \code{date_to} \\
\route{/events/<id>}                        & GET      & all     & \code{events.event_detail} \\
\route{/events/<id>/register}               & POST     & login   & \code{events.register_for_event} \\
\route{/events/<id>/cancel}                 & POST     & login   & \code{events.cancel_registration}; form field \code{return_to=my} returns to My registrations \\
\route{/my-registrations}                   & GET      & login   & \code{events.my_registrations} \\
\midrule
\route{/admin/}                             & GET      & admin   & \code{admin.dashboard} \\
\route{/admin/venues}                       & GET      & admin   & \code{admin.list_venues} \\
\route{/admin/venues/new}                   & GET/POST & admin   & \code{admin.create_venue} \\
\route{/admin/venues/<id>/edit}             & GET/POST & admin   & \code{admin.edit_venue} \\
\route{/admin/venues/<id>/delete}           & POST     & admin   & \code{admin.delete_venue} \\
\route{/admin/events}                       & GET      & admin   & \code{admin.list_events}: all events incl.\ past \\
\route{/admin/events/new}                   & GET/POST & admin   & \code{admin.create_event} \\
\route{/admin/events/<id>/edit}             & GET/POST & admin   & \code{admin.edit_event} \\
\route{/admin/events/<id>/delete}           & POST     & admin   & \code{admin.delete_event} \\
\route{/admin/events/<id>/participants}     & GET      & admin   & \code{admin.participants} \\
\route{/admin/events/<id>/}\newline\hspace*{1em}\route{participants.csv} & GET & admin & \code{admin.participants_csv}: CSV download \\
\route{/admin/reports}                      & GET      & admin   & \code{admin.reports_page} \\
\midrule
\route{/static/<path>}                      & GET      & all     & \code{static}: \code{style.css} \\
\bottomrule
\caption{All 22 URL rules (from \texttt{flask --app app routes}). \emph{login}: visitors are
redirected to the login page. \emph{admin}: visitors are redirected, students get 403.
\texttt{<id>} is an integer; anything else gives 404.}
\end{xltabular}}

% =====================================================================
\section{Module reference}

Signatures are shown in grey boxes. ``Returns'' describes the value; routes return a rendered page
or a redirect.

\subsection{\texorpdfstring{\code{app/db.py}}{app/db.py}: database access}

\fn{get_db() -> sqlite3.Connection}{Returns the connection for the current request, opening it on
first use and storing it in \code{g.db}. Sets \code{row_factory = sqlite3.Row} and
\code{PRAGMA foreign_keys = ON}.}

\fn{close_db(error=None)}{Closes \code{g.db} if one was opened. Registered with
\code{app.teardown_appcontext}, so it runs after every request and app context.}

\fn{init_db()}{Runs \code{sql/schema.sql} with \code{executescript}. Destroys existing data.}

\fn{add_indexes()}{Runs \code{sql/indexes.sql}. Idempotent (\code{IF NOT EXISTS}).}

\fn{init_db_command / add_indexes_command / create_admin_command}{Click commands \code{init-db},
\code{add-indexes}, \code{create-admin --name --email --password}. \code{create-admin} skips an
existing email and stores \code{generate_password_hash(password)}.}

\fn{init_app(app)}{Registers \code{close_db} and the three commands on the app.}

\subsection{\texorpdfstring{\code{app/auth.py}}{app/auth.py}: accounts and access control}

\fn{class User(UserMixin)}{Built from a \code{users} row; attributes \code{id}, \code{name},
\code{email}, \code{role}; property \code{is_admin}. \code{UserMixin} supplies
\code{is_authenticated}, \code{get_id()} and friends.}

\fn{load_user(user_id) -> User | None}{The Flask-Login user loader: one \code{SELECT} by id on every
request that uses \code{current_user}. Returns \code{None} for a deleted user, which logs them out.}

\fn{admin_required(view)}{Decorator. Not logged in: \code{login_manager.unauthorized()} (redirect to
login). Logged in but not admin: \code{abort(403)}. Uses \code{functools.wraps} so Flask still sees the
original function name as the endpoint.}

\fn{register()}{GET shows the form; POST validates name (non-empty), email (contains \code{@}, lower-cased)
and password (at least 6 characters), then inserts with role \code{'student'}. A duplicate email
raises \code{sqlite3.IntegrityError}, shown as a friendly message. Logged-in users are redirected.}

\fn{login()}{Looks up the email, verifies with \code{check_password_hash}, calls
\code{login_user(User(row))}. The same message is used for an unknown email and a wrong password.}

\fn{logout()}{\code{@login_required}; calls \code{logout_user()} and redirects home.}

\subsection{\texorpdfstring{\code{app/events.py}}{app/events.py}: browsing and registrations}

\fn{CATEGORIES = ["Technical", "Cultural", "Sports", "Workshop", "Seminar"]}{The single list of
categories, imported by \code{admin.py} and \code{seed.py}.}

\fn{now_string() -> str}{Current local time as \code{'YYYY-MM-DD HH:MM'}.}

\fn{get_event(event_id) -> Row | None}{One event joined with its venue and organiser, plus
\code{seats_left}. \code{LEFT JOIN registrations ... GROUP BY e.id}.}

\fn{is_registered(user_id, event_id) -> bool}{\code{SELECT 1} on the unique pair.}

\fn{build_search_sql(keyword="", category="", venue_id=None, date_from="", date_to="") -> (str, list)}{
Builds the search query. Starts with \code{e.start_time >= now}; adds a condition for each non-empty
filter. \code{conditions} only ever contains fixed SQL written in the code; user input goes into
\code{params}. Keyword matches title or description with \texttt{LIKE '\%text\%'} (case-insensitive for
ASCII). \code{date_to} has \code{" 23:59"} appended so the whole day is included. Separate from execution
so \code{profiling/explain.py} can explain exactly the same SQL.}

\fn{search_events(...) -> list[Row]}{Runs \code{build_search_sql} with the same arguments.}

\fn{list_events()}{Reads filters from \code{request.args} (\code{venue_id} with \code{type=int}, so bad
input becomes \code{None}), calls \code{search_events}, loads venues for the dropdown and passes the
filters back to keep the form filled in.}

\fn{event_detail(event_id)}{404 if missing; computes \code{is_registered} and \code{has_started} for
the template's action area.}

\fn{register_for_event(event_id)}{POST, login required. Applies the registration rules
(Section~\ref{sec:rules}) and inserts; \code{registered_at} comes from the column default.}

\fn{cancel_registration(event_id)}{POST, login required. Refuses after the start; otherwise
\code{DELETE} and uses \code{cursor.rowcount} to tell whether a row existed. Redirects to My
registrations when \code{return_to=my}, otherwise to the detail page.}

\fn{my_registrations()}{The current user's upcoming registrations joined with event and venue,
ordered by start time.}

\subsection{\texorpdfstring{\code{app/admin.py}}{app/admin.py}: administration}

All routes are decorated with \code{@admin_required}.

\fn{to_db_time(value) -> str | None}{\code{'2026-11-20T14:30'} $\rightarrow$
\code{'2026-11-20 14:30'}; \code{None} if it does not parse.}

\fn{find_clash(venue_id, start_time, end_time, ignore_event_id=None) -> Row | None}{First event at the
venue whose interval overlaps: \code{start_time < new_end AND end_time > new_start}. Touching
intervals do not clash. \code{ignore_event_id} excludes the event being edited.}

\fn{read_event_form(form, event_id=None) -> (dict, None) | (None, str)}{Validates and normalises the
event form (Section~\ref{sec:rules}). Returns the clean values or the first error message.}

\fn{read_venue_form(form, venue_id=None) -> (dict, None) | (None, str)}{Name and location required,
capacity a positive integer; when editing, not below \code{MAX(capacity)} of the venue's events.}

\fn{get_participants(event_id) -> list[Row]}{Name, email, \code{registered_at}, earliest first.}

\fn{dashboard()}{One query with five scalar sub-queries: events, upcoming, venues, students,
registrations.}

\fn{list_venues(), create_venue(), edit_venue(venue_id), delete_venue(venue_id)}{Venue CRUD. Forms
redisplay the submitted values on error. Delete refuses while events use the venue (the foreign key
would also block it).}

\fn{list_events(), create_event(), edit_event(event_id), delete_event(event_id)}{Event CRUD.
\code{created_by} is \code{current_user.id}. Edit pre-fills \code{datetime-local} values by replacing the
space with \code{T}. Delete relies on \code{ON DELETE CASCADE} to remove registrations and returns 404 when
nothing was deleted.}

\fn{all_venues(), get_event_or_404(event_id)}{Small helpers for the event form dropdown and the
participant pages.}

\fn{participants(event_id), participants_csv(event_id)}{HTML list, and the same rows written with
\code{csv.writer} into \code{io.StringIO}, returned as \code{text/csv} with
\code{Content-Disposition: attachment; filename=event_<id>_participants.csv}.}

\fn{reports_page()}{Renders \code{admin/reports.html} with the three report functions.}

\subsection{\texorpdfstring{\code{app/reports.py}}{app/reports.py}: report queries}

The SQL is held in module constants so \code{profiling/explain.py} can explain the same text.

\fn{percent(part, whole) -> float}{\code{round(100 * part / whole, 1)}, or 0.0 when \code{whole} is 0.
Done in Python because Python and SQLite round halves differently, and the numbers had to match the
first version exactly.}

\fn{event_stats() -> list[dict]}{\code{EVENT_STATS_SQL}: every event with \code{registered}
(\code{LEFT JOIN} + \code{COUNT(r.id)}), sorted by registrations, start, id; adds \code{fill_rate}.}

\fn{popular_categories() -> list[dict]}{\code{CATEGORIES_SQL}: per category,
\code{COUNT(DISTINCT e.id)} events (the join repeats events) and \code{COUNT(r.id)} registrations.}

\fn{venue_utilisation() -> list[dict]}{\code{VENUES_SQL}: a derived table counts registrations per
event first, then the outer query sums events, registrations and capacity per venue (so capacity is
not multiplied by the number of registrations). \code{LEFT JOIN} keeps venues without events.
Adds \code{utilisation}.}

\begin{lstlisting}[language=SQL, caption={VENUES\_SQL: the most involved report query.}]
SELECT v.id, v.name, v.location,
       COUNT(e.id) AS events,
       COALESCE(SUM(e.registered), 0) AS registrations,
       COALESCE(SUM(e.capacity), 0)   AS seats_offered
FROM venues v
LEFT JOIN (
    SELECT ev.id, ev.venue_id, ev.capacity, COUNT(r.id) AS registered
    FROM events ev
    LEFT JOIN registrations r ON r.event_id = ev.id
    GROUP BY ev.id
) AS e ON e.venue_id = v.id
GROUP BY v.id
ORDER BY registrations DESC, v.name
\end{lstlisting}

\subsection{\texorpdfstring{\code{app/seed.py}}{app/seed.py}: demo data}

\fn{seed_command(num_events=50, num_students=100)}{The \code{flask seed} command. Seeds \code{random}
and Faker with 42 for reproducible data. Requires an admin (the organiser). Deletes registrations,
events, venues and students, resets their \code{sqlite_sequence} counters, inserts the 10 fixed
\code{VENUES}, students with one shared pre-computed password hash (hashing is deliberately slow),
events on random non-clashing slots in the next 60 days (gives up after \code{20 * N} attempts), and
for each event a \code{random.sample} of distinct students no larger than its capacity.}

\fn{overlaps(start, end, booked_slots) -> bool}{The same overlap rule as \code{find_clash}, on
Python \code{datetime} pairs.}

% =====================================================================
\section{Business rules}
\label{sec:rules}

\subsection{Registration}

\code{register_for_event} checks, in this order, and stops at the first failure:

\begin{enumerate}
  \item user is not an admin;
  \item \code{start_time > now} (the event has not started);
  \item no existing registration (also guaranteed by \code{UNIQUE (user_id, event_id)});
  \item \code{seats_left > 0}.
\end{enumerate}

Cancellation is allowed only before the event starts.

\subsection{Event validation (\texorpdfstring{\code{read_event_form}}{read\_event\_form})}

\begin{enumerate}
  \item title not empty;
  \item category in \code{CATEGORIES};
  \item start and end are valid \code{datetime-local} values;
  \item end after start (also a \code{CHECK} constraint);
  \item capacity is an integer above 0 (also a \code{CHECK});
  \item venue exists, and capacity $\le$ venue capacity;
  \item when editing: capacity $\ge$ current registrations;
  \item no clash at the venue (\code{find_clash}, ignoring the event itself when editing).
\end{enumerate}

\subsection{The overlap rule}

Two half-open intervals $[s_1, e_1)$ and $[s_2, e_2)$ overlap exactly when
$s_1 < e_2 \land e_1 > s_2$. This single condition covers partial overlap on either side,
containment in both directions and identical intervals, while letting back-to-back events share a
boundary.

\subsection{Venue rules}

Name and location required; capacity a positive integer; on edit, not smaller than the largest event
at the venue; delete only when no events use the venue.

% =====================================================================
\section{Templates and styling}

\subsection{Template inheritance}

Every page extends \code{templates/base.html}, which defines the \code{title} and \code{content}
blocks, the navigation bar (links depend on \code{current_user}), the flash-message area and the
footer. Jinja2 auto-escapes all variables, so user-supplied text such as event descriptions cannot
inject HTML.

\begin{table}[H]
\centering
\small
\begin{tabularx}{\linewidth}{@{}l X@{}}
\toprule
\textbf{Template} & \textbf{Rendered by / purpose} \\
\midrule
\code{base.html}                    & layout shared by all pages \\
\code{index.html}                   & \code{index}: welcome text and counts \\
\code{error.html}                   & 403 and 404 handlers (\code{code}, \code{message}) \\
\code{auth/login.html}, \code{auth/register.html} & login and sign-up forms \\
\code{events/list.html}             & search form and event cards \\
\code{events/detail.html}           & one event and the Register / Cancel action area \\
\code{events/my_registrations.html} & student's upcoming registrations with Cancel buttons \\
\code{admin/dashboard.html}         & counters and shortcuts \\
\code{admin/venues.html}, \code{admin/venue_form.html} & venue list; shared new/edit form (\code{venue_id} is \code{None} for new) \\
\code{admin/events.html}, \code{admin/event_form.html} & event list; shared new/edit form \\
\code{admin/participants.html}      & participant table and Export CSV button \\
\code{admin/reports.html}           & three report tables and CSS fill-rate bars \\
\bottomrule
\end{tabularx}
\caption{The 15 templates.}
\end{table}

\begin{warning}[Jinja comments]
Jinja processes tags even inside HTML comments, so \verb|<!-- {% block %} -->| breaks the template.
Use Jinja's own comment syntax \verb|{# ... #}| for comments that mention template tags.
\end{warning}

\subsection{CSS}

\code{static/style.css} is the only stylesheet. Colours and sizes are CSS custom properties on
\code{:root} (\texttt{-{}-primary: \#1d4ed8}, \code{--error}, \code{--radius} \dots), so the palette is
changed in one place. Layout uses flexbox (header, sticky footer, search form that wraps) and CSS
grid (\code{repeat(auto-fill, minmax(280px, 1fr))} for event cards, so the number of columns follows
the screen width). Wide tables sit in a wrapper with \code{overflow-x: auto} for phones. Fill-rate bars
on the reports page are a \code{span} inside a \code{div} whose \code{width} style is the percentage.

% =====================================================================
\section{Security model}

\begin{table}[H]
\centering
\small
\begin{tabularx}{\linewidth}{@{}l X@{}}
\toprule
\textbf{Threat} & \textbf{Mitigation} \\
\midrule
SQL injection        & every value through \code{?} placeholders; SQL text never built from input \\
HTML injection (XSS) & Jinja2 auto-escaping \\
Stolen passwords     & salted, slow hashes (Werkzeug); generic login error message \\
Privilege escalation & role fixed to \code{'student'} at sign-up; \code{admin_required} on every admin route; DB \code{CHECK} on role \\
Cross-site requests  & all state changes are POST; \code{SameSite=Lax} session cookie \\
Accidental GET side effects & delete/register/cancel are forms, never links \\
Bad IDs              & \code{<int:...>} URL converters and 404s \\
\bottomrule
\end{tabularx}
\caption{Security measures.}
\end{table}

\paragraph{Known limitations.} No CSRF tokens (\code{SameSite} covers most cases); the seat check and
the insert are two statements, so two simultaneous requests for the last seat could over-book by one;
\code{/auth/logout} is a GET; the default \code{SECRET_KEY} must be replaced in production; keyword
search cannot use an index because of the leading \texttt{\%} wildcard.

% =====================================================================
\section{Testing}

Tests live in \code{tests/} and run with \code{./scripts/test.sh}. \code{pytest.ini} sets
\code{testpaths = tests} and \code{pythonpath = .} so tests can \code{import app}.

\begin{table}[H]
\centering
\small
\begin{tabularx}{\linewidth}{@{}l X@{}}
\toprule
\textbf{Fixture} & \textbf{Provides} \\
\midrule
\code{app}            & \code{create_app} with \code{TESTING=True} and a database in \code{tmp_path}; schema applied \\
\code{client}         & \code{app.test_client()}: a fake browser \\
\code{auth}           & \code{AuthActions} with \code{register()} and \code{login()} helpers \\
\code{admin_user}     & an admin inserted directly with SQL \\
\code{sample_events}  & one venue and four events (three upcoming, one past); returns \texttt{\{title: id\}} \\
\code{student_client} & client logged in as a newly registered student \\
\code{admin_client}   & client logged in as the admin \\
\bottomrule
\end{tabularx}
\caption{Fixtures in \texttt{tests/conftest.py}.}
\end{table}

When you add a feature, add tests next to the related ones and run the whole suite. The Test Plan \&
Report lists every existing test case.

% =====================================================================
\section{Profiling toolkit}

\begin{description}
  \item[\code{scripts/profile.sh <label> [N]}] creates \code{profiling/profile.db} (via
        \code{CAMPUS_DB}), seeds N events, optionally adds indexes, then runs the two modules below.
  \item[\code{profiling/run_profile.py}] for each workload in \code{WORKLOADS}: \code{RUNS} timed runs
        with \code{time.perf_counter()} and SQL statements counted with
        \code{conn.set_trace_callback}; then one run under \code{cProfile}. Writes
        \code{profile.prof}, \code{profile_stats.txt} (top 20 by cumulative time) and
        \code{timings.txt}.
  \item[\code{profiling/explain.py}] writes \code{EXPLAIN QUERY PLAN} output for the search variants,
        the naive per-event count, the three report queries, the clash check and the
        ``is registered'' query to \code{query_plans.txt}.
\end{description}

To profile a new function, add a \code{("name", function)} pair to \code{WORKLOADS}, and its SQL to
\code{queries()} in \code{explain.py}. View a profile with \code{snakeviz profiling/<label>/profile.prof}.
The saved results in \code{profiling/before/} are the evidence of the naive version and must not be
overwritten; results and analysis are in \code{profiling/RESULTS.md} and the project report.

% =====================================================================
\section{Extending the application}

\subsection{Add a page}
\begin{enumerate}
  \item Add a view function to the right Blueprint, with \code{@login_required} or
        \code{@admin_required} as needed.
  \item Put any logic worth testing in a plain function next to it.
  \item Create a template whose first line is \verb|{% extends "base.html" %}|, and link
        to the page with \code{url_for("blueprint.function")}.
  \item Add tests for the happy path, each refusal, and access control.
\end{enumerate}

\subsection{Change the schema}
\begin{enumerate}
  \item Edit \code{sql/schema.sql} (constraints belong in the database where possible).
  \item Update every \code{INSERT}/\code{SELECT} that touches the table, the seed command and fixtures.
  \item Recreate the database with \code{./scripts/setup.sh} (there are no migrations; this wipes data).
  \item If a new query filters or joins on a column, check its plan with \code{EXPLAIN QUERY PLAN} and
        add an index to \code{sql/indexes.sql} only if it helps.
\end{enumerate}

\subsection{Add a category or a report}
Categories: add the name to \code{CATEGORIES} in \code{events.py} (and optionally to \code{TITLES} in
\code{seed.py}). Reports: write one SQL query as a constant in \code{reports.py}, a function that runs
it, a table in \code{admin/reports.html}, and a test with exact expected numbers.

\subsection{Coding conventions}
\begin{itemize}
  \item PEP 8 style, descriptive names, a docstring on every function, comments that explain why.
  \item SQL keywords in upper case; one clause per line in longer queries; table aliases
        \code{e}, \code{v}, \code{r}, \code{u}.
  \item Every data-changing request is POST and ends in a redirect.
  \item Bash scripts: \code{set -euo pipefail}, source \code{common.sh}, clear \code{echo} progress
        messages, LF line endings (enforced by \code{.gitattributes}).
\end{itemize}

\subsection{Git workflow}
Small commits on \code{main}, each leaving the tests passing, using Conventional Commit prefixes:
\code{feat:}, \code{fix:}, \code{test:}, \code{perf:}, \code{chore:}, \code{docs:}. Generated files
(\code{venv/}, \code{*.db}, \code{__pycache__}, LaTeX build files) are in \code{.gitignore}.

\end{document}
